% GENERATED FILE. Edit canonical lessons, then run npm run generate:books.
% canonical-source-hash: fnv1a64:a90508cbd39c7e6b
% canonical-lessons: FR-C26-entendre, FR-C26-dormir, FR-C26-marcher, FR-C26-courir

\chapter{Hearing, Sleeping, Walking, Running}
\label{ch:fr-hearing-sleeping-walking-running}
\hlchaptermodality{26}

\begin{chapteropening}
\textbf{By the end of this chapter:} \emph{I can say that I hear, sleep, walk and run in French, use ça marche for \textquotedblleft{}that works,\textquotedblright{} and say why entendre carries English's intend inside it while the Latin verb for hearing ended up in English instead.}

Run the chapter's four verbs together with je — j'entends, je dors, je marche, je cours — then use courir's spelling to settle why canis became chien while currere stayed hard at the front: the ca- shift only ever struck c before a. Reaches back past the chapter to le chien and its palatalization from Chapter 22, le chat from the same chapter, il fait chaud and il pleut from Chapter 21, and the ne … pas negation from Chapter 18.
\end{chapteropening}

\section[entendre]{entendre — \textquotedblleft{}to hear,\textquotedblright{} and the stretching inside it}

\label{lesson:FR-C26-entendre}



\begin{quote}
You can think and understand. Now the ear — and French's word for it is English's \emph{intend}.
\end{quote}

\subsection*{You'll want to know: the word}
\begin{quote}
\textbf{entendre} — to hear. Two nasal \emph{an} sounds in a row, then the \emph{-dre} of \emph{prendre}: say it \emph{ahn-TAHN-druh}.
\end{quote}

>

\begin{quote}
\textbf{J'entends le chien.} — I hear the dog.
\end{quote}

The \emph{-ds} of \emph{j'entends} is silent, as in \emph{je prends}.

\begin{grammarlens}[title={Grammar Lens: the nasal that stays put}]
\emph{Prendre}'s nasal dies in the plural. \emph{Entendre} looks like its twin and does the opposite — the stem never moves:

\begin{itemize}
\raggedright
  \item \textbf{j'entends} — \emph{jahn-TAHN}
  \item \textbf{tu entends} — \emph{ahn-TAHN}
  \item \textbf{il entend} — \emph{ahn-TAHN} (three identical to the ear)
  \item \textbf{nous entendons} — \emph{ahn-tahn-DOHN}, nasal still there
  \item \textbf{vous entendez} — \emph{ahn-tahn-DAY}
  \item \textbf{ils entendent} — \emph{ahn-TAHND}
\end{itemize}
\end{grammarlens}

\begin{cousinweb}
\textbf{entendre} $\leftarrow$ Latin \textbf{intendere} — \textbf{in-} \textquotedblleft{}toward\textquotedblright{} plus \textbf{tendere}, \textquotedblleft{}\textbf{to stretch}.\textquotedblright{} To \emph{intendere} was to stretch a thing out and aim it: a bowstring, a journey, your attention. English kept the family — \textbf{intend}, plainest of them, plus \textbf{intent}, \textbf{intense}, \textbf{tension}, a \textbf{tent} (a stretched skin), \textbf{extend} and \textbf{pretend}.

Set your three mind-verbs together: French pictures each as a different thing the body does.

\begin{itemize}
\raggedright
  \item \textbf{comprendre} — to \textbf{grasp together}, built on \emph{prendre}.
  \item \textbf{penser} — to \textbf{weigh}, from \emph{pēnsāre}.
  \item \textbf{entendre} — to \textbf{stretch toward}.
\end{itemize}
\end{cousinweb}

\begin{culture}
From its oldest French records \emph{entendre} already covers a span at once: to hear, to attend, to understand, to intend. Nothing arrived late. What happened since is that the other senses fell away and hearing took the word outright — because Latin's verb for hearing, \textbf{audīre}, wore down in French to \emph{ouïr} and stopped being usable. English kept that family instead: \textbf{audio}, \textbf{audience}, \textbf{audible}, \textbf{audition}.

The abandoned senses survive in phrases French never rebuilt:

\begin{itemize}
\raggedright
  \item \textbf{entendu} — \textquotedblleft{}understood,\textquotedblright{} \textquotedblleft{}agreed.\textquotedblright{}
  \item \textbf{s'entendre} — to hear one another, and so to \textbf{get along}.
  \item \textbf{l'entente} — an understanding between parties, borrowed into English whole in \emph{Entente Cordiale}.
\end{itemize}

So the two languages split one Latin pair: French hears with the verb English uses for intending, and English hears with the verb French let go.
\end{culture}

\subsection*{Your turn}
Say these aloud:

\begin{itemize}
\raggedright
  \item \textquotedblleft{}entendre\textquotedblright{} — \emph{ahn-TAHN-druh}, two nasals
  \item \textquotedblleft{}j'entends le chien\textquotedblright{}
  \item \textquotedblleft{}nous prenons\textquotedblright{} then \textquotedblleft{}nous entendons\textquotedblright{} — one flattens, one holds
  \item \textquotedblleft{}j'écris, je comprends, j'entends\textquotedblright{}
  \item the three pictures — stretching, grasping, weighing
  \item \textquotedblleft{}intend, tension, tent\textquotedblright{} — all a stretching
\end{itemize}

\begin{tcolorbox}[breakable,colback=teal!4,colframe=teal!35!black,title={Before you move on}]
Say \textquotedblleft{}I hear the dog\textquotedblright{} and \textquotedblleft{}we hear.\textquotedblright{} (\emph{J'entends le chien. Nous entendons} — and unlike \emph{nous prenons}, nothing flattens.) What did \emph{intendere} mean, and which English verb is its plainest survival? (\textbf{To stretch toward}; \textbf{intend}.) Why did French hand hearing to this verb? (Latin's \emph{audīre} wore down to \emph{ouïr} and fell out of use — English kept the \emph{audīre} words instead.) What does \emph{s'entendre} mean? (\textbf{To get along}.)
\end{tcolorbox}
\section[dormir]{dormir — \textquotedblleft{}to sleep,\textquotedblright{} and the room named after it}

\label{lesson:FR-C26-dormir}



\begin{quote}
After a verb that changed jobs, here is one that never changed. \emph{Dormir} is Latin's word for sleeping, handed on almost untouched — and it hides one trap.
\end{quote}

\subsection*{You'll want to know: the word}
\begin{quote}
\textbf{dormir} — to sleep. Said \emph{dor-MEER}, with the throaty French \textbf{r} in the middle and at the end.
\end{quote}

>

\begin{quote}
\textbf{Le chat dort.} — The cat is sleeping.
\end{quote}

French has no separate \textquotedblleft{}is sleeping\textquotedblright{} form. \emph{Il dort} covers both \textquotedblleft{}he sleeps\textquotedblright{} and \textquotedblleft{}he is sleeping\textquotedblright{}; the plain present does all the work.

\begin{grammarlens}[title={Grammar Lens: the m that goes missing}]
The stem is \emph{dorm-}, and in the singular French throws the \textbf{m} away. Listen for where it comes back:

\begin{itemize}
\raggedright
  \item \textbf{je dors} — \emph{dor}
  \item \textbf{tu dors} — \emph{dor}
  \item \textbf{il dort} — \emph{dor} (three identical, the \emph{m} gone)
  \item \textbf{nous dormons} — \emph{dor-MOHN}, and there it is
  \item \textbf{vous dormez} — \emph{dor-MAY}
  \item \textbf{ils dorment} — \emph{DORM}
\end{itemize}

One rule, and you have the whole verb: singular loses the last stem consonant, plural gets it back. A small family of \emph{-ir} verbs shares the shape, so learning it once is worth more than one verb.
\end{grammarlens}

\begin{cousinweb}
\textbf{dormir} $\leftarrow$ Latin \textbf{dormīre}, \textquotedblleft{}to sleep.\textquotedblright{} Spanish kept \emph{dormir}, Italian \emph{dormire}, Portuguese \emph{dormir} — the verb crossed into its daughter languages untouched, which is rarer than it sounds. English has no inherited descendant; it borrowed the family later, through French:

\begin{itemize}
\raggedright
  \item \textbf{dormitory} $\leftarrow$ \emph{dormītōrium}, a sleeping place; French wore the identical word down to \textbf{dortoir}.
  \item \textbf{dormant} — Old French \emph{dormant}, \textquotedblleft{}sleeping\textquotedblright{}; a dormant volcano is asleep.
  \item \textbf{dormer} — a roof window, from Old French \emph{dormeor}, because the room behind it was where you slept.
\end{itemize}

And an honest stop. The \textbf{dormouse} looks like the best cousin of the lot and is not one. The story needs an Anglo-Norman word for \textquotedblleft{}sleeper\textquotedblright{} that nobody has found in a text, and \emph{dormouse}'s real origin is unknown. A guess in circulation, not an etymology.
\end{cousinweb}

\subsection*{Your turn}
Say these aloud:

\begin{itemize}
\raggedright
  \item \textquotedblleft{}dormir\textquotedblright{} — \emph{dor-MEER}, throaty \emph{r}
  \item \textquotedblleft{}je dors\textquotedblright{} then \textquotedblleft{}nous dormons\textquotedblright{} — the \emph{m} vanishes, then returns
  \item \textquotedblleft{}le chat dort\textquotedblright{} — the cat, the travelling word out of Egypt
  \item \textquotedblleft{}j'entends le chat\textquotedblright{} then \textquotedblleft{}le chat dort\textquotedblright{}
  \item \textquotedblleft{}dormir, dormitory, dormant\textquotedblright{} — one Latin verb, borrowed back
  \item \textquotedblleft{}je comprends\textquotedblright{} then \textquotedblleft{}je dors\textquotedblright{}
\end{itemize}

\begin{tcolorbox}[breakable,colback=teal!4,colframe=teal!35!black,title={Before you move on}]
Say \textquotedblleft{}the cat is sleeping\textquotedblright{} and \textquotedblleft{}we sleep.\textquotedblright{} (\emph{Le chat dort. Nous dormons} — the \emph{m} is silent in the singular and audible in the plural.) Where does \emph{dormitory} come from? (Latin \emph{dormītōrium}, \textquotedblleft{}a \textbf{sleeping place}\textquotedblright{} — and French wore the same word down to \emph{dortoir}.) Is the \textbf{dormouse} in this family? (\textbf{No} — that is folk etymology; the word's real origin is unknown.) And \emph{chat} — an inherited Latin word, or a traveller? (\textbf{A traveller}, from \emph{cattus}, moving with the cat itself out of Egypt.)
\end{tcolorbox}
\section[marcher]{marcher — \textquotedblleft{}to walk,\textquotedblright{} and the step hiding inside French negation}

\label{lesson:FR-C26-marcher}



\begin{quote}
You can hear and you can sleep. Now move. This verb does a second job too, used dozens of times a day.
\end{quote}

\subsection*{You'll want to know: the word}
\begin{quote}
\textbf{marcher} — to walk. Said \emph{mar-SHAY}: the \emph{-er} is a clean \emph{ay}, and the \emph{ch} is \emph{sh}, never \emph{k}.
\end{quote}

>

\begin{quote}
\textbf{Je marche.} — I walk.
\end{quote}

After \emph{dormir}'s vanishing \emph{m}, this one asks nothing: an ordinary \textbf{-er} verb. \emph{Je marche, tu marches, il marche, nous marchons, vous marchez, ils marchent} — four of the six sound identical, just \emph{marsh}.

\begin{culture}
\emph{Marcher} is also the everyday verb for a thing \textbf{working}. A machine, a plan, an arrangement: if it functions, it walks.

\begin{quote}
\textbf{Ça marche.} — \textquotedblleft{}That works.\textquotedblright{} Also \textquotedblleft{}Fine,\textquotedblright{} \textquotedblleft{}Deal,\textquotedblright{} \textquotedblleft{}Sounds good.\textquotedblright{} \textbf{Ça ne marche pas.} — \textquotedblleft{}That doesn't work.\textquotedblright{}
\end{quote}

There is your \emph{ne … pas} from the answer-word \emph{non}, round a real verb at last. And \emph{ça marche} alone is an agreement, two words long.

\emph{Entendre} once held several jobs and shed all but hearing. \emph{Marcher} holds two right now.
\end{culture}

\begin{cousinweb}
Here the honest answer is that nobody is certain. Old French \emph{marchier} first meant \textquotedblleft{}\textbf{to trample, to tread on},\textquotedblright{} and two accounts compete:

\begin{itemize}
\raggedright
  \item Frankish \textbf{markōn}, \textquotedblleft{}to mark, to press a step in\textquotedblright{} — preferred by most reference works.
  \item Late Latin \emph{marcāre}, from \textbf{marcus}, \textquotedblleft{}\textbf{a hammer}\textquotedblright{} — a blow and a tread being one motion.
\end{itemize}

Three consequences, in falling order of certainty:

\begin{itemize}
\raggedright
  \item English \textbf{march}, the verb, is borrowed straight \textbf{from} \emph{marcher}. Secure whichever origin is right.
  \item English \textbf{mark} is Germanic, from Old English \emph{mearc}. A relative \textbf{only if} the Frankish account holds — conditional, not a fact.
  \item French \textbf{la marche}, a borderland, with \textbf{marquis} and German \textbf{margrave}, comes from the Germanic \textbf{noun} \emph{marka}, \textquotedblleft{}boundary\textquotedblright{} — a separate arrival.
\end{itemize}

And one cousin you own already. The \textbf{pas} in \emph{ne … pas} began as Latin \emph{passus}, \textquotedblleft{}\textbf{a step}.\textquotedblright{} \emph{Je ne marche pas} is, underneath, \textquotedblleft{}I do not walk a step.\textquotedblright{}
\end{cousinweb}

\subsection*{Your turn}
Say these aloud:

\begin{itemize}
\raggedright
  \item \textquotedblleft{}marcher\textquotedblright{} — \emph{mar-SHAY}, \emph{sh} not \emph{k}
  \item \textquotedblleft{}je marche\textquotedblright{} then \textquotedblleft{}nous marchons\textquotedblright{}
  \item \textquotedblleft{}ça marche\textquotedblright{} — that works, that's fine
  \item \textquotedblleft{}non\textquotedblright{}, then \textquotedblleft{}je ne marche pas\textquotedblright{} — and \emph{pas} is a \textbf{step}
  \item \textquotedblleft{}j'entends, je dors, je marche\textquotedblright{}
  \item \textquotedblleft{}nous dormons, nous marchons\textquotedblright{} — one stem shifts, one never does
\end{itemize}

\begin{tcolorbox}[breakable,colback=teal!4,colframe=teal!35!black,title={Before you move on}]
Say \textquotedblleft{}I walk\textquotedblright{} and \textquotedblleft{}that works.\textquotedblright{} (\emph{Je marche. Ça marche}.) What are the two competing origins? (Frankish \textbf{markōn\emph{, \textquotedblleft{}to press a step in,\textquotedblright{} or Latin }marcus\emph{, \textquotedblleft{}\textbf{a hammer}.\textquotedblright{} Genuinely unsettled.) Which English word is certainly related, and which only might be? (\textbf{March}, borrowed straight from }marcher\emph{; \textbf{mark} only if the Frankish account is right.) What did }pas\emph{ mean before it meant \textquotedblleft{}not\textquotedblright{}? (\textbf{A step}.)}}
\end{tcolorbox}
\section[courir]{courir — \textquotedblleft{}to run,\textquotedblright{} and why it is not \emph{chourir}}

\label{lesson:FR-C26-courir}



\begin{quote}
Hearing, sleeping, walking — and now the fastest of the four. Its Latin root is one of the busiest in English, and its spelling settles a question the dog raised.
\end{quote}

\subsection*{You'll want to know: the word}
\begin{quote}
\textbf{courir} — to run. Said \emph{koo-REER}, throaty at both \emph{r}s.
\end{quote}

>

\begin{quote}
\textbf{Le chien court.} — The dog runs.
\end{quote}

\begin{grammarlens}[title={Grammar Lens: the same endings, a different stem}]
\emph{Courir} takes \emph{dormir}'s singular endings — \textbf{-s, -s, -t}, all silent:

\begin{itemize}
\raggedright
  \item \textbf{je cours} — \emph{koor}
  \item \textbf{tu cours} — \emph{koor}
  \item \textbf{il court} — \emph{koor}
  \item \textbf{nous courons} — \emph{koo-ROHN}
  \item \textbf{vous courez} — \emph{koo-RAY}
  \item \textbf{ils courent} — \emph{KOOR}
\end{itemize}

But the resemblance stops there. \emph{Dormir} \textbf{sheds} a consonant in the singular: \emph{dorm-} becomes \emph{dor-}. \emph{Courir} sheds nothing — its stem \emph{cour-} has nothing spare to lose.
\end{grammarlens}

\begin{cousinweb}
\emph{Marcher} left a choice between a Frankish verb and a Latin hammer. This one leaves none. \textbf{courir} $\leftarrow$ Latin \textbf{currere}, \textquotedblleft{}to run,\textquotedblright{} which English mined harder than almost any other verb:

\begin{itemize}
\raggedright
  \item \textbf{current} and \textbf{currency} both run; a \textbf{cursor} runs across a screen; a \textbf{course} is a run, a \textbf{corridor} a running-place.
  \item \textbf{courier}, one who runs; \textbf{curriculum}, literally \textquotedblleft{}\textbf{a running}.\textquotedblright{}
  \item \textbf{concur}, run together; \textbf{occur}, run up against; \textbf{excursion}.
\end{itemize}

Now the spelling question. \emph{Chien} came from \emph{canis}, and its \textbf{c} turned into \textbf{ch} — the shift that made \emph{champ} and \emph{chanter}. So why is this verb not \emph{chourir}? Because that shift only struck \textbf{c} before \textbf{a}. \emph{Currere} had \textbf{c} before \textbf{u}, untouched, so the \emph{k} survived. One rule, two regular outcomes.

One honest limit. Latin \textbf{carrus}, \textquotedblleft{}wagon,\textquotedblright{} gave English \textbf{car} and \textbf{cargo}, and descends from the same very old root. But Latin borrowed \emph{carrus} from Gaulish: a cousin by another road, not a word built out of \emph{currere}.
\end{cousinweb}

\subsection*{Your turn}
Say these aloud:

\begin{itemize}
\raggedright
  \item the chapter's four with \emph{je} — \textquotedblleft{}j'entends, je dors, je marche, je cours\textquotedblright{}
  \item \textquotedblleft{}le chien court\textquotedblright{} then \textquotedblleft{}le chat dort\textquotedblright{}
  \item \textquotedblleft{}il fait chaud, je cours\textquotedblright{}
  \item \textquotedblleft{}il pleut, je ne cours pas\textquotedblright{}
  \item \textquotedblleft{}je dors, nous dormons\textquotedblright{} then \textquotedblleft{}je cours, nous courons\textquotedblright{} — one sheds
  \item \textquotedblleft{}chien\textquotedblright{} then \textquotedblleft{}courir\textquotedblright{} — \emph{ch} before \emph{a}, \emph{k} before \emph{u}
  \item \textquotedblleft{}current, course, courier, curriculum\textquotedblright{} — every one a running
  \item \textquotedblleft{}ça marche\textquotedblright{} — the agreement, two words
\end{itemize}

\begin{tcolorbox}[breakable,colback=teal!4,colframe=teal!35!black,title={Before you move on}]
Say the chapter's four verbs with \emph{je}. (\emph{J'entends, je dors, je marche, je cours}.) What does \textbf{curriculum} literally mean? (\textbf{A running} — from \emph{currere}.) Why did \emph{canis} become \emph{chien} while \emph{currere} stayed hard? (The shift only hit \textbf{c} before \textbf{a}.) Which of the four began as a \textbf{stretching}? (\emph{\textbf{Entendre}} — \emph{intendere}, to stretch toward.) Now: \textquotedblleft{}the dog runs\textquotedblright{} and \textquotedblleft{}it is hot.\textquotedblright{} (\emph{Le chien court. Il fait chaud}.)
\end{tcolorbox}
